\section{Discussion and Final Remarks}
\label{sec:final}

We investigated the fragment $\xpath$ of the query language XPath from a DL perspective. 
This fragment features data comparisons and is interpreted over data graphs with functional accessibility relations. We showed with an example that $\xpath$ is expressive enough to model non-trivial scenarios.

With this as a starting point, we presented a novel connection between $\xpath$ and a fragment of $\alcd$. This connection enabled us to view $\xpath$ as a DL with a very simple concrete domain and to obtain upper bounds on the complexity of several reasoning tasks via transfer results. We then improved on these bounds and showed that concept satisfiability, ABox consistency, and concept satisfiability w.r.t.\ a TBox are all \np-complete for $\xpath$. This is in stark contrast to more expressive DLs, such as those involving arithmetical domains, for which concept satisfiability is $\nexp$-hard~\cite{Lutz01,Lutz02b,Lutz02}.

Our results highlight the relevance of the restrictions we imposed on the expressive power of $\xpath$. In particular, the requirement that accessibility relations be partial functions plays a crucial role in our selection function, which establishes a small model property for the reasoning tasks under investigation.
Beyond the functional restriction, simple tests, and other comparisons in the concrete domain, may likewise lead to interesting complexity bounds for fragments of XPath over arbitrary accessibility relations

As we discussed in the introduction, only the purely navigational fragment of XPath has been studied in a DL setting before~\cite{KostylevRV14,Kost:xpat15}. In this article, we show how
to deal with data comparisons by linking XPath with DL with concrete domains.  The work we presented here is a first step in this direction. In particular, it would be interesting to see whether our methods and results can be adapted and/or extended to deal with other logics, such as those capturing fragments of the Shapes Constraint Language (SHACL) (see, e.g.,~\cite{Ortiz23,Ahmetaj0S23} for details), or more interestingly, to establish complexity forfragment of Graph Query Language (GQL) (see, e.g.,~\cite{FrancisGGLMMMPR23,FrancisGGLMMMPR23b,FigueiraLP26}).  We should also investigate the effect 
of allowing additional constructs on path expressions like tests, and non-deterministic choice. Finally, we would also like to implement $\xpath$ reasoners for the tasks we investigated.